\section{Introduction}\label{sec:intro}

A US trademark application carries a description of the goods or services
the mark will cover. That description is written to secure scope of
protection, it is dated, it is public, and it is filed by the millions of
firms that never patent. The literature that reads trademarks as
innovation data has worked mostly from what surrounds the description:
how many marks a firm files, in which classes, how broadly, and whether
it renews them \citep{Castaldi2020, Flikkema2019, Dinlersoz2018}. The
text itself has entered only piecemeal, as breadth counts or similarity
scores. This paper builds two research assets from the description and
the prosecution record around it, states what each can and cannot
support, and makes both public.

The first asset is a corpus. The USPTO publishes every application since
1884 with its full prosecution history, but the literature has mostly
read one field from it: the terminal status code. Status codes run
together events that the maintenance schedule keeps apart. Between the
fifth and sixth year after registration, an owner must file a sworn
declaration that the mark is still in use in commerce, and must renew
again at year ten; a registration whose owner does neither is cancelled,
and both cancellations share a status code. We re-parse all 13.99 million
case files into 242 million dated events, so each cancellation carries
its date and the record shows which deadline the mark failed. The
five-year proof of continued use thereby becomes a dated, sworn,
product-level survival record, available for every registered mark more
than seven years old. The same parse records office actions and
responses, counsel of record, filing basis and owner domicile, and
owners are linked by normalized name to SEC reporting, Regulation~D
rounds, and PatentsView patent assignees. Where the existing public
datasets supply case-file snapshots \citep{Graham2013} or Census-linked
registration histories \citep{Dinlersoz2018}, the event-dated corpus
adds the full prosecution timeline of each case, and with it the dated
outcomes.

The second asset is a measure. A topic model fitted once over the whole
corpus reduces each description to a mix of themes --- whether the
filing describes an online offering, say, or an artificial-intelligence
application --- and we compare each filing's theme mix, by
Kullback--Leibler divergence, with what its own Nice class filed in the
five years before its date and in the five years after. Their average is
the filing's \emph{atypicality}: how unusual its language is for its
industry. Their difference is its \emph{lead}: whether its language sat
ahead of or behind where the industry's language was moving. Looking
both backward and forward is what separates \emph{unusual} from
\emph{early}, and the two components predict different things.

We validate the measure by variation rather than by a single fit. The
identical filings are rescored under three partitions of the vocabulary,
two reference weightings, five window constructions, a second model seed,
and with and without the 47\% of registrations whose text is shared
verbatim with another filing. Lead survives every variation and predicts
failure at the five-year proof: $+1.4$ points from the most lagging to
the most leading fifth with class-by-cohort fixed effects, drafting and
filer controls and owner-clustered errors ($t = 8.3$); $+2.3$, $+2.6$
and $+1.9$ points raw under the three partitions. Atypicality does not
survive: its apparently protective effect at the proof falls from $4.4$
points under global themes to $0.7$ under per-class themes, and its sign
at the SEC-reporting margin flips with how description length is held.

Two hazards follow for any research that scores trademark text, and we
document both on identical samples. Scoring words rather than theme
distributions measures description length: term-scored atypicality
correlates with log distinct terms at $-0.68$. And firm-level
correlations built on word scoring --- the kind the literature reports
--- dissolve or reverse under theme scoring on the same firms; what they
carry is unsigned lexical specificity, not vocabulary position.

What the measure predicts about pioneering and industry waves is the
subject of a companion paper; here those outcomes appear only as
validation, and the paper makes no claim about mechanism.
Section~\ref{sec:record} describes the corpus and the outcomes it dates.
Section~\ref{sec:measure} builds the measure. Section~\ref{sec:related}
settles nomenclature against the innovation literature's terms.
Section~\ref{sec:validation} tests the measure against registration, the
five-year proof, patenting and SEC reporting.
Section~\ref{sec:robust} summarizes what changes under every alternative
construction tried, and Section~\ref{sec:practices} states the practices
that follow for trademark-text research.
